\documentclass[12pt]{article}
\usepackage[a4paper,margin=1in]{geometry}
\usepackage[T1]{fontenc}
\usepackage{mathptmx}
\usepackage{amsmath,amssymb,setspace,enumitem,booktabs,tabularx,titlesec}
\usepackage{pdfpages}
\usepackage[hidelinks]{hyperref}
\doublespacing
\setlength{\parindent}{0pt}
\setlength{\parskip}{1pt}
\titlespacing*{\section}{0pt}{5pt}{2pt}
\titlespacing*{\subsection}{0pt}{4pt}{1pt}
\setlength{\abovedisplayskip}{5pt}
\setlength{\belowdisplayskip}{5pt}
\setlength{\abovedisplayshortskip}{3pt}
\setlength{\belowdisplayshortskip}{3pt}
\setlist{nosep,leftmargin=1.5em}
\newcommand{\E}{\mathcal E}
\newcommand{\pr}{\operatorname{pr}}
\newcommand{\eps}{\varepsilon}

\begin{document}

\includepdf[pages=1,fitpaper=true]{Sample+First_Page.pdf}

\begin{center}
\textbf{Lie-Symmetry Reduction and Semi-Analytical Neural Approximation\\
of Parameter-Dependent Nonlinear PDEs}\\[4pt]
\end{center}

\section{Research Problem and Motivation}

This study aims to reduce complex partial differential equations (PDEs) into simpler forms and obtain exact or approximate analytical solutions, with applications in mathematical physics and engineering. Lie point symmetries identify continuous transformations that map solutions to other solutions. For suitable invariant families, these transformations combine space and time into one variable, reducing the PDE to an ordinary differential equation (ODE).

The study joins this reduction with machine learning: analytical calculations determine the known solution terms, and a neural network learns the remaining correction in the reduced variables. The guiding example is the published Saint--Venant--Exner (SVE) master equation near the symmetry-enhanced exponent $m=3$ \cite{lagrange}.

Paliathanasis, Leon, and Leach begin with the one-dimensional Eulerian
Saint--Venant--Exner system
\begin{align}
h_t+(hu)_x&=0, \nonumber\\
(hu)_t+\left(hu^2+\tfrac12gh^2\right)_x+ghB_x&=0, \label{eq:eulerian}\\
B_t+\xi Q_x&=0,\qquad Q=A_g u^m,\qquad
\beta=mA_g\xi, \nonumber
\end{align}
where $h$ is water depth, $u$ is depth-averaged velocity, $B$ is bed
elevation, $g$ is gravitational acceleration, and $\xi=(1-\epsilon_p)^{-1}$ accounts for bed porosity $\epsilon_p$. With the particle position $x=\phi(t,s)$,
\[
\phi_t=u(t,\phi),\qquad \phi_s=h(t,\phi)^{-1},
\qquad \widetilde B(t,s)=B(t,\phi(t,s)),
\]
the source paper obtains the scalar master equation
\begin{equation}
\E_m[\phi]:=
\partial_t\left(\phi_s\phi_{tt}-g\frac{\phi_{ss}}{\phi_s^2}\right)
-g\beta\partial_s\left(\phi_t^{m-1}\frac{\phi_{ts}}{\phi_s}\right)=0.
\label{eq:master}
\end{equation}
The model requires one qualification. Applying the chain rule to the Eulerian bed equation gives
\begin{equation}
\widetilde B_t-\frac{\phi_t}{\phi_s}\widetilde B_s
+\beta\phi_t^{m-1}\frac{\phi_{ts}}{\phi_s}=0.
\label{eq:lag-exner-audit}
\end{equation}
The transport term $-\phi_t\widetilde B_s/\phi_s$ is absent from the paper's displayed Lagrangian bed equation. This study therefore treats \eqref{eq:master} as the published mathematical benchmark; equivalence to the physical SVE system must be checked before interpreting its solutions as water and sediment flows.

For arbitrary $m$, its admitted generators include
\[
X_1=\partial_t,\quad X_2=\partial_s,\quad X_3=\partial_\phi,\quad
X_4=t\partial_t+s\partial_s+\phi\partial_\phi.
\]
At $m=3$, the algebra gains
\[
X_5=t\partial_t-2s\partial_s.
\]
The finite symmetry maps will be stated explicitly as
\[
S_i(\eps):=S_{X_i}(\eps)=\exp(\eps X_i).
\]
In particular,
\[
S_5(\eps):(t,s,\phi)\mapsto(e^\eps t,e^{-2\eps}s,\phi),
\]
which leaves $\sigma=st^2$ invariant and gives $\phi(t,s)=F(\sigma)$. The additional symmetry is generally lost away from $m=3$. We therefore study $m=3+\delta$ to determine how the analytical reduction changes with the transport exponent.
After nondimensionalization with fixed reference scales, $A_g$ and $\xi$ are held fixed, so $\beta(\delta)=(3+\delta)A_g\xi$.

The choice is geometric, not universal: the source allows $1\le m\le4$, but
the algebra of \eqref{eq:master} increases in dimension precisely at $m=3$.
Continuation will determine the useful $\delta$ range; none is assumed.

Symmetry-informed PINNs and symbolic Taylor terms with neural remainders already exist \cite{liepinn,taylor}. The proposed focus is their connection through approximate symmetry reduction near a parameter value where the symmetry algebra enlarges.

\section{Research Question and Core Proposition}

Can symmetry reduction near $m=3$ yield exact or approximate analytical profiles, and can a learned remainder improve those profiles while retaining a clear error estimate?

The analytical target is a conditional error statement. Let $Y$ solve the derived first-order ODE system and let $\widehat Y$ share its initial data. On an interval $I$, suppose both remain in a regular region where the ODE is Lipschitz with constant $L$. If the uniform ODE residual of $\widehat Y$ is at most $\rho$, then
\begin{equation}
\|\widehat Y-Y\|_{\infty}
\le \frac{\rho}{L}\left(e^{L|I|}-1\right),
\label{eq:profile-bound}
\end{equation}
with $\rho|I|$ when $L=0$. For an order-$p$ approximate reduction, we will derive conditions under which the reconstructed PDE residual satisfies
\begin{equation}
\|\E_{3+\delta}[\widehat\phi_\theta]\|_\infty
\le C_{\rm lift}\rho+C_{\rm sym}|\delta|^{p+1}.
\label{eq:lift-bound}
\end{equation}
Here $\widehat\phi_\theta$ is the reconstructed approximation. The two terms separate the reduced residual from the truncation of the parameter expansion. The constants must be bounded on the selected domain; sampled loss alone is insufficient. Neither estimate asserts an error bound against an unrestricted PDE solution.

\section{Objectives}

\begin{enumerate}
\item Determine relevant generators $X_i$ and finite maps $S_i(\eps)$, and derive their invariant reductions.
\item Obtain exact profiles where possible, and construct series or perturbation solutions otherwise.
\item Learn the unresolved reduced correction while preserving analytical terms and prescribed data.
\item Establish conditions relating reduced residuals to profile error and reconstructed PDE residual.
\end{enumerate}

\section{Methodology}

\subsection{Complete Lie-symmetry analysis}

Start from the general point generator
\[
X=\tau(t,s,\phi)\partial_t+\zeta(t,s,\phi)\partial_s+
\eta(t,s,\phi)\partial_\phi.
\]
Extend $X$ to act on derivatives through third order \cite{olver} and solve
\[
\left.\pr^{(3)}X(\E_m)\right|_{\E_m=0}=0
\]
by matching coefficients of the independent derivatives. This gives equations for $\tau,\zeta,\eta$ at generic and exceptional $m$. Commutators check the algebra, and adjoint transformations identify equivalent reductions. Integrating each generator gives its finite map $S_i(\eps)$. Each selected reduction will include the equations defining its invariants.

For $X_5$,
\[
\frac{dt}{d\eps}=t,\qquad \frac{ds}{d\eps}=-2s,\qquad
\frac{d\phi}{d\eps}=0
\]
gives $\sigma=st^2$ and $\phi=F(\sigma)$.  Substitution at $m=3$ yields
\[
\E_3[F(st^2)]=t^{-1}R_3[F],\qquad
R_3[F]=a(\sigma,F')F'''+b(\sigma,F',F''),
\]
where $a$ and $b$ follow directly from differentiation.
On $r=F'>0$ and $a\ne0$, set $v=F''$ to obtain $Y'=(r,v,-b/a)$ for $Y=(F,r,v)$. First integrals and further symmetries will be used to seek exact solutions; otherwise the ODE determines Taylor coefficients recursively. Only data compatible with the invariant family can be imposed.

\subsection{Approximate symmetry near \texorpdfstring{$m=3$}{m=3}}

Under $X_5$, the two terms in \eqref{eq:master} have scaling weights $-1$ and $2-m$. Their mismatch explains why the extra symmetry is lost when $m\ne3$. Following approximate symmetry theory \cite{approx}, a corrected generator
\[
X_\delta=X_5+\delta X^{(1)}+\cdots+\delta^pX^{(p)}
\]
will be sought by expanding the equation and invariance condition order by order. For positive dimensionless velocity,
\[
\phi_t^{2+\delta}=\phi_t^2
\left[1+\delta\log\phi_t+O(\delta^2)\right]
\]
Since $\phi_t=2\sigma F'/t$, a candidate correction must allow logarithmic terms:
\[
\phi^{[p]}(t,s;\delta)=
\sum_{k=0}^p\delta^k\sum_{\ell=0}^k(\log t)^\ell F_{k\ell}(st^2).
\]
Coefficient matching tests whether this ansatz closes into ODEs for $F_{k\ell}$; closure is not assumed. The expansion uses compact regions with $t>0$, $\phi_s>0$, and $\phi_t>0$, where logarithmic corrections remain small. If the corrected symmetry fails to exist, the remaining PDE residual identifies the limit of the approximation.

\subsection{Semi-analytic neural approximation}

For reduced data $F(\sigma_0)=F_*$, $F'(\sigma_0)=r_*$, and $F''(\sigma_0)=v_*$, use the exact-data trial
\[
\widehat F_\theta(\sigma)=F_*+r_*z+\tfrac12v_*z^2+z^3N_\theta(\widehat\sigma,\delta),
\qquad z=\sigma-\sigma_0.
\]
Here $\widehat\sigma$ is a scaled coordinate. The factor $z^3$ prevents the learned remainder from changing the three prescribed initial values. For a Taylor polynomial of degree $q\ge2$, use $z^{q+1}$ instead. Boundary-value problems require a factor preserving their boundary data.

Training minimizes the squared reduced-equation residual at sampled profile and parameter values. At $m=3$, this is $\widehat F_\theta'''+b/a$; away from three, the derived correction ODEs replace it. Symbolic algebra checks the reduction, and automatic differentiation evaluates the neural residual. Network size and sampling are selected by convergence, without arbitrary fixed values. The reconstructed solution is then substituted back into the PDE to measure the residual left by the complete approximation.

\section{Expected Outcomes}

The study will produce explicit symmetry transformations $S_i(\eps)$, reduced ODEs, and exact or controlled approximate solution families for the selected benchmark. The semi-analytical neural representation will preserve known coefficients and data while approximating unresolved profiles. Its error analysis will identify where the approximation remains reliable and where symmetry breaking or singularities limit its use.

The broader contribution is a reusable procedure: identify a useful symmetry, derive a simpler equation, solve as much as possible analytically, and learn the remaining correction. This provides a basis for parameterized wave and transport problems in mathematical physics and engineering, with physical applications conditional on verifying the model and its data.

\begin{singlespace}\small
\begin{thebibliography}{9}
\bibitem{lagrange}
A.~Paliathanasis, G.~Leon, and P.~G.~L.~Leach,
``Lie symmetries and similarity solutions for a shallow-water model with bed elevation in Lagrange variables,''
\emph{Mathematics}, 14, 433, 2026.

\bibitem{approx}
V.~A.~Baikov, R.~K.~Gazizov, and N.~Kh.~Ibragimov,
``Approximate symmetries,'' \emph{Math. USSR-Sbornik}, 64(2), 427--441, 1989.

\bibitem{liepinn}
T.~Akhound-Sadegh et al.,
``Lie point symmetry and physics-informed networks,''
\emph{Advances in Neural Information Processing Systems}, 36, 2023.

\bibitem{taylor}
C.~K.~Nagesh et al.,
``Taylor-model physics-informed neural networks for ordinary differential equations,''
International Conference on Neuro-symbolic Systems, 2025.

\bibitem{olver}
P.~J.~Olver, \emph{Applications of Lie Groups to Differential Equations}, 2nd ed., Springer, 1993.
\end{thebibliography}
\end{singlespace}

\end{document}
